\chapter{Seleção de Instruções}\label{chap:x86-codegen}

\section{Introdução}

Com a IRT canonicalizada e o \texttt{Coloring} em mãos, a fase final do
back-end é a \textbf{seleção de instruções}: traduzir cada nó da IRT para
uma sequência concreta de instruções X86-64. Diferentemente das fases
anteriores, que manipulam abstrações (temporários, grafos, conjuntos),
a seleção de instruções lida com \emph{registradores físicos, offsets de
memória e o protocolo de chamada de função}.

Dois desafios técnicos merecem atenção especial nesta fase:

\begin{enumerate}
  \item \textbf{Avaliação de \texttt{BINOP}}: compilar $e_1 \oplus e_2$
    requer guardar o resultado de $e_1$ enquanto $e_2$ é avaliado. Se
    $e_2$ contém uma chamada de função, todos os registradores caller-saved
    são clobberados. A solução é empilhar $e_1$ na pilha e contabilizar
    esse deslocamento extra para manter o alinhamento de 16 bytes nas
    chamadas aninhadas.
  \item \textbf{Alinhamento de pilha}: cada \texttt{callq} exige
    \texttt{\%rsp} múltiplo de 16. O gerador deve rastrear dinamicamente
    quantos bytes extras foram empilhados (além do prólogo) para inserir
    padding quando necessário.
\end{enumerate}

\section{Notação Formal}

Usamos dois julgamentos para especificar a tradução:

\begin{itemize}
  \item $\Phi \vdash e \Rightarrow \vec{I}$: sob o colorindo $\Phi$
    (que mapeia cada temporário a um registrador físico ou offset de memória),
    a expressão $e$ é compilada para a sequência de instruções $\vec{I}$.
    O resultado fica sempre em \texttt{\%rax} (\textit{scratchAcc}).

  \item $\Phi \vdash s \Rightarrow \vec{I}$: o enunciado $s$ é compilado
    para $\vec{I}$ sem valor de retorno explícito.
\end{itemize}

Denotamos:
$r_{\mathit{acc}} = \texttt{\%rax}$,
$r_{\mathit{bin}} = \texttt{\%r11}$,
$r_{\mathit{mem}} = \texttt{\%r10}$.

\section{Regras de Compilação de Expressões}

\subsection{Constante e Temporário}

\[
  \infer[\text{CE-Const}]
  {\Phi \vdash \texttt{CONST}(n) \Rightarrow [\texttt{movq } \$n,\,
  \%\mathit{rax}]}{}
\]

\[
  \infer[\text{CE-Temp-Reg}]
  {\Phi \vdash \texttt{TEMP}(t) \Rightarrow [\texttt{movq } r,\,
  \%\mathit{rax}]}
  {\Phi(t) = r \quad r \neq r_{\mathit{acc}}}
  \qquad
  \infer[\text{CE-Temp-Acc}]
  {\Phi \vdash \texttt{TEMP}(t) \Rightarrow []}
  {\Phi(t) = r_{\mathit{acc}}}
\]

\[
  \infer[\text{CE-Temp-Mem}]
  {\Phi \vdash \texttt{TEMP}(t) \Rightarrow [\texttt{movq }
  \mathit{off}(\%\mathit{rbp}),\, \%\mathit{rax}]}
  {\Phi(t) = \mathit{off}}
\]

\subsection{Operação Binária}

\[
  \infer[\text{CE-Binop}]
  {\Phi \vdash \texttt{BINOP}(\oplus, e_1, e_2) \Rightarrow
    \vec{I_1};\; \texttt{pushq \%rax};\; \vec{I_2};\;
  \texttt{popq \%r11};\; \llbracket\oplus\rrbracket}
  {\Phi \vdash e_1 \Rightarrow \vec{I_1} \quad \Phi \vdash e_2
  \Rightarrow \vec{I_2}}
\]

O \texttt{pushq \%rax} salva $e_1$ na pilha \emph{antes} de avaliar $e_2$.
Isso protege o valor de $e_1$ mesmo que $e_2$ contenha chamadas que
clobberem registradores caller-saved. O \texttt{popq \%r11} restaura $e_1$
no registrador scratch de BINOP, deixando \texttt{\%rax = $e_2$} e
\texttt{\%r11 = $e_1$}.

A tradução $\llbracket\oplus\rrbracket$ de cada operador binário (com
\texttt{\%r11 = $e_1$} e \texttt{\%rax = $e_2$}) é:

\begin{center}
  \begin{tabular}{ll}
    \toprule
    Operador & Instruções $\llbracket\oplus\rrbracket$ \\
    \midrule
    \texttt{BAdd} & \texttt{addq \%r11, \%rax} \\
    \texttt{BMul} & \texttt{imulq \%r11, \%rax} \\
    \texttt{BSub} & \texttt{negq \%rax; addq \%r11, \%rax} \\
    \texttt{BDiv} & \texttt{movq \%rax, \%r10; movq \%r11, \%rax;
    cqo; idivq \%r10} \\
    \texttt{BMod} & (idem BDiv) \texttt{movq \%rdx, \%rax} \\
    \texttt{BAnd} & \texttt{andq \%r11, \%rax} \\
    \texttt{BOr}  & \texttt{orq \%r11, \%rax} \\
    \texttt{BXor} & \texttt{xorq \%r11, \%rax} \\
    \texttt{BEq}  & \texttt{cmpq \%rax, \%r11; sete \%al; movzbq \%al, \%rax} \\
    \texttt{BNe}  & \texttt{cmpq \%rax, \%r11; setne \%al; movzbq
    \%al, \%rax} \\
    \texttt{BLt}  & \texttt{cmpq \%rax, \%r11; setl \%al; movzbq \%al, \%rax} \\
    \texttt{BLe}  & \texttt{cmpq \%rax, \%r11; setle \%al; movzbq
    \%al, \%rax} \\
    \texttt{BGt}  & \texttt{cmpq \%rax, \%r11; setg \%al; movzbq \%al, \%rax} \\
    \texttt{BGe}  & \texttt{cmpq \%rax, \%r11; setge \%al; movzbq
    \%al, \%rax} \\
    \bottomrule
  \end{tabular}
\end{center}

Para comparações, \texttt{cmpq src, dst} calcula $\mathit{dst} - \mathit{src}$
e atualiza as flags. Portanto \texttt{cmpq \%rax, \%r11} calcula
$\texttt{\%r11} - \texttt{\%rax} = e_1 - e_2$, que é exatamente o que
as condições de comparação exigem.

Para subtração, usamos a identidade $e_1 - e_2 = -(e_2) + e_1 = e_1 + (-e_2)$:
\texttt{negq} nega $e_2$ em \texttt{\%rax} e depois \texttt{addq} soma $e_1$
de \texttt{\%r11}.

Para divisão e módulo, o X86-64 exige que o dividendo esteja em
\texttt{rdx:rax} (estendido por \texttt{cqo}). Usamos \texttt{\%r10} como
scratch para guardar o divisor antes de mover o dividendo para \texttt{\%rax}.

\subsection{Memória e ESEQ}

\[
  \infer[\text{CE-Mem}]
  {\Phi \vdash \texttt{MEM}(e) \Rightarrow \vec{I};\; [\texttt{movq }
  0(\%\mathit{rax}),\, \%\mathit{rax}]}
  {\Phi \vdash e \Rightarrow \vec{I}}
\]

\[
  \infer[\text{CE-ESeq}]
  {\Phi \vdash \texttt{ESEQ}(s, e) \Rightarrow \vec{I_s};\; \vec{I_e}}
  {\Phi \vdash s \Rightarrow \vec{I_s} \quad \Phi \vdash e
  \Rightarrow \vec{I_e}}
\]

A regra CE-ESeq é fornecida por completude. Na prática, a canonicalização
elimina todos os nós \texttt{ESEQ} antes desta fase.

\subsection{Chamada de Função}

Seja $\vec{r} = [r_1, \ldots, r_n]$ os registradores de argumento da ABI
(\texttt{\%rdi}, \texttt{\%rsi}, \texttt{\%rdx}, \texttt{\%rcx},
\texttt{\%r8}, \texttt{\%r9}):

\[
  \infer[\text{CE-Call}]
  {\Phi \vdash \texttt{CALL}(\texttt{NAME}(f), [a_1, \ldots, a_n])
    \Rightarrow
    \mathit{save};\;
    \vec{I_1};\; [\texttt{movq \%rax,} r_1];\;
    \cdots;\;
    \vec{I_n};\; [\texttt{movq \%rax,} r_n];\;
    [\texttt{xorq \%rax, \%rax};\; \texttt{callq } f];\;
  \mathit{restore}}
  {\Phi \vdash a_i \Rightarrow \vec{I_i}\; \forall i}
\]

Onde $\mathit{save}$ e $\mathit{restore}$ salvam e restauram os
registradores caller-saved do pool de alocação que estão em uso (i.e.,
  atribuídos a algum temporário pelo coloridor e portanto com valores que
devem sobreviver ao \texttt{callq}). O \texttt{xorq \%rax, \%rax} zera
\texttt{\%eax} como exigido pelo ABI para chamadas variádicas
(\texttt{printf}, etc.).

\section{Regras de Compilação de Enunciados}

\subsection{MOVE}

\[
  \infer[\text{CS-Move-Reg}]
  {\Phi \vdash \texttt{MOVE}(\texttt{TEMP}(t), e) \Rightarrow
  \vec{I};\; [\texttt{movq \%rax,} r]}
  {\Phi \vdash e \Rightarrow \vec{I} \quad \Phi(t) = r \quad r \neq
  r_{\mathit{acc}}}
\]

\[
  \infer[\text{CS-Move-Acc}]
  {\Phi \vdash \texttt{MOVE}(\texttt{TEMP}(t), e) \Rightarrow \vec{I}}
  {\Phi \vdash e \Rightarrow \vec{I} \quad \Phi(t) = r_{\mathit{acc}}}
\]

\[
  \infer[\text{CS-Move-Spill}]
  {\Phi \vdash \texttt{MOVE}(\texttt{TEMP}(t), e) \Rightarrow \vec{I};\;
  [\texttt{movq \%rax,} \mathit{off}(\%\mathit{rbp})]}
  {\Phi \vdash e \Rightarrow \vec{I} \quad \Phi(t) = \mathit{off}}
\]

\[
  \infer[\text{CS-Move-Mem}]
  {\Phi \vdash \texttt{MOVE}(\texttt{MEM}(\mathit{addr}),
    \mathit{val}) \Rightarrow
    \vec{I_a};\; [\texttt{movq \%rax, \%r10}];\; \vec{I_v};\;
  [\texttt{movq \%rax,} 0(\%r10)]}
  {\Phi \vdash \mathit{addr} \Rightarrow \vec{I_a} \quad \Phi \vdash
  \mathit{val} \Rightarrow \vec{I_v}}
\]

Na regra CS-Move-Mem, \texttt{\%r10} (\textit{scratchMem}) guarda o
endereço enquanto $\mathit{val}$ é avaliado em \texttt{\%rax}. Isso
funciona porque \texttt{\%r10} não pertence ao pool de alocação e,
portanto, nunca contém o valor de um temporário.

\subsection{Rótulos, Saltos e Retorno}

\[
  \infer[\text{CS-Label}]
  {\Phi \vdash \texttt{LABEL}(l) \Rightarrow [l\texttt{:}]}{}
  \qquad
  \infer[\text{CS-Jump}]
  {\Phi \vdash \texttt{JUMP}(\texttt{NAME}(l)) \Rightarrow [\texttt{jmp } l]}{}
\]

\[
  \infer[\text{CS-CJump}]
  {\Phi \vdash \texttt{CJUMP}(e, l_t, l_f) \Rightarrow
    \vec{I};\; [\texttt{testq \%rax, \%rax};\; \texttt{jne } l_t;\;
  \texttt{jmp } l_f]}
  {\Phi \vdash e \Rightarrow \vec{I}}
\]

\[
  \infer[\text{CS-Return-Void}]
  {\Phi \vdash \texttt{RETURN}([]) \Rightarrow \mathit{epilogue}(\Phi)}{}
\]

\[
  \infer[\text{CS-Return-Val}]
  {\Phi \vdash \texttt{RETURN}([e]) \Rightarrow \vec{I};\;
  \mathit{epilogue}(\Phi)}
  {\Phi \vdash e \Rightarrow \vec{I}}
\]

O epílogo $\mathit{epilogue}(\Phi)$ executa: restaura os registradores
callee-saved em ordem inversa ao prólogo, restaura \texttt{\%rbp} e executa
\texttt{retq}. O resultado de retorno fica em \texttt{\%rax}, que já contém
o valor de $e$ após a regra CE acima.

\section{Alinhamento de Pilha}

\subsection{O Problema}

A instrução \texttt{callq} exige \texttt{\%rsp} alinhado a 16 bytes.
O prólogo garante o alinhamento \emph{estático} usando \texttt{frameAdj}.
Mas durante a execução do corpo, \texttt{pushq \%rax} do BINOP
(regra CE-Binop) empilha 8 bytes extras, quebrando o alinhamento para
qualquer \texttt{callq} dentro de $e_2$.

\subsection{Monada de Geração com Rastreamento de Profundidade}

O gerador de código rastreia a profundidade extra acumulada pelos
\texttt{pushq} de BINOP numa mônada de estado:

\begin{minted}{haskell}
data MState = MState [Instr] Int
--                           ↑ extraDepth

type X86M = State MState

pushScratch :: Reg -> X86M ()
pushScratch r = do
  emit (Pushq (Reg r))
  modify (\(MState is d) -> MState is (d + 8))

popScratch :: Reg -> X86M ()
popScratch r = do
  emit (Popq (Reg r))
  modify (\(MState is d) -> MState is (d - 8))

getExtraDepth :: X86M Int
getExtraDepth = gets (\(MState _ d) -> d)
\end{minted}

O campo \texttt{extraDepth} conta apenas os bytes de \texttt{pushScratch}
(BINOPs), não os do prólogo (já contabilizados em \texttt{frameAdj}).

\subsection{Cálculo de Padding}

Antes de cada \texttt{callq}, o gerador:

\begin{enumerate}
  \item Empilha os registradores caller-saved em uso (\texttt{saved},
    $n_{\mathit{saved}}$ registradores, $8 n_{\mathit{saved}}$ bytes).
  \item Verifica se o deslocamento total $(\mathit{extra} + 8
    n_{\mathit{saved}})$
    é múltiplo de 16. Se não, empilha 8 bytes de padding.
\end{enumerate}

\begin{minted}{haskell}
compileCall :: Env -> String -> [Expr] -> X86M ()
compileCall env f args = do
  extra <- getExtraDepth
  let saved   = envCSRegs env
      nSaved  = length saved
      needPad = (extra + 8 * nSaved) `mod` 16 /= 0
  mapM_ (emit . Pushq . Reg) saved
  when needPad $ emit (Subq (Imm 8) (Reg RSP))
  mapM_ (uncurry loadArg) (zip args argRegs)
  emit (Xorq (Reg RAX) (Reg RAX))
  emit (Callq f)
  when needPad $ emit (Addq (Imm 8) (Reg RSP))
  mapM_ (emit . Popq . Reg) (reverse saved)
  where
    loadArg e r = do
      compileExpr env e
      when (r /= scratchAcc) $ emit (Movq (Reg scratchAcc) (Reg r))
\end{minted}

\begin{remark}
  O padding usa \texttt{subq \$8, \%rsp} em vez de \texttt{pushq} para
  não empilhar lixo interpretável como dado. O ajuste contrário
  \texttt{addq \$8, \%rsp} é feito antes de restaurar os caller-saved.
\end{remark}

\section{Implementação: Módulo \texttt{IRToX86}}

\subsection{Ambiente de Compilação}

\begin{minted}{haskell}
data Env = Env
  { envCol    :: Coloring   -- temporário → registrador ou slot
  , envCSaved :: [Reg]      -- callee-saved usados (para o epílogo)
  , envFAdj   :: Int        -- F do subq no prólogo (0 = omitido)
  , envCSRegs :: [Reg]      -- caller-saved do pool em uso
  }
\end{minted}

\subsection{Prólogo e Epílogo}

\begin{minted}{haskell}
prologue :: [Reg] -> Int -> [Instr]
prologue csSaved fAdj =
  [ Pushq (Reg RBP)
  , Movq  (Reg RSP) (Reg RBP)
  ]
  ++ map (Pushq . Reg) csSaved
  ++ [Subq (Imm fAdj) (Reg RSP) | fAdj > 0]

epilogue :: Env -> [Instr]
epilogue env =
  [Addq (Imm (envFAdj env)) (Reg RSP) | envFAdj env > 0]
  ++ map (Popq . Reg) (reverse (envCSaved env))
  ++ [Popq (Reg RBP), Retq]
\end{minted}

Os registradores callee-saved são desempilhados em \emph{ordem inversa}:
a pilha é LIFO.

\subsection{Carregamento de Parâmetros}

Após o prólogo, os parâmetros chegam nos registradores da ABI
(\texttt{\%rdi}, \texttt{\%rsi}, \ldots). Se o coloridor atribuiu a um
parâmetro $t$ um registrador diferente do ABI, é necessário um
\texttt{movq}:

\begin{minted}{haskell}
loadParams :: [Temp] -> Coloring -> [Instr]
loadParams params col = concat $ zipWith load params argRegs
  where
    load t r = case Map.lookup t col of
      Just (InReg dr) | dr /= r -> [Movq (Reg r) (Reg dr)]
      Just (InReg _)             -> []
      Just (InMem off)           -> [Movq (Reg r) (Mem off RBP)]
      Nothing                    -> []
\end{minted}

\subsection{Compilação de Expressões}

\begin{minted}{haskell}
compileExpr :: Env -> Expr -> X86M ()
compileExpr _   (CONST n)  = emit (Movq (Imm n) (Reg scratchAcc))
compileExpr env (TEMP t)   = loadTemp t (envCol env)
compileExpr env (BINOP op e1 e2) = do
  compileExpr env e1
  pushScratch scratchAcc
  compileExpr env e2
  popScratch  scratchBinop
  compileBinOp op
compileExpr env (MEM e) = do
  compileExpr env e
  emit (Movq (Mem 0 scratchAcc) (Reg scratchAcc))
compileExpr env (CALL (NAME f) args) = compileCall env f args
compileExpr env (ESEQ s e) = do
  compileStmt env s
  compileExpr env e
\end{minted}

\subsection{Compilação de Enunciados}

\begin{minted}{haskell}
compileStmt :: Env -> Stmt -> X86M ()
compileStmt env (MOVE (TEMP t) e) = do
  compileExpr env e
  storeTemp t (envCol env)
compileStmt env (MOVE (MEM addrE) valE) = do
  compileExpr env addrE
  emit (Movq (Reg scratchAcc) (Reg scratchMem))
  compileExpr env valE
  emit (Movq (Reg scratchAcc) (Mem 0 scratchMem))
compileStmt _   (LABEL l)           = emit (LocLabel l)
compileStmt _   (JUMP (NAME l))     = emit (Jmp l)
compileStmt env (CJUMP cond lt lf)  = do
  compileExpr env cond
  emit (Testq (Reg scratchAcc) (Reg scratchAcc))
  emit (Jcc NE lt)
  emit (Jmp lf)
compileStmt env (RETURN [])    = mapM_ emit (epilogue env)
compileStmt env (RETURN (e:_)) = do
  compileExpr env e
  mapM_ emit (epilogue env)
compileStmt env (SEQ s1 s2) = do
  compileStmt env s1
  compileStmt env s2
\end{minted}

\subsection{Compilação de Funções}

\begin{minted}{haskell}
compileFunc :: FuncDef -> [Instr]
compileFunc fd =
  let body0             = canonicalize (funcBody fd)
      (col, spillBytes) = colorFunc (funcParams fd) body0
      csSaved           = calleeSavedInCol col
      fAdj              = frameAdj (length csSaved) spillBytes
      csRegs            = callerSavedInCol col
      env               = Env col csSaved fAdj csRegs
      stmts             = linearize body0
      body              = run (mapM_ (compileStmt env) stmts)
  in Dir "" :
     Dir (".globl " ++ funcName fd) :
     GlobLabel (funcName fd) :
     prologue csSaved fAdj ++
     loadParams (funcParams fd) col ++
     body ++
     epilogue env
\end{minted}

O epílogo ao final é um \textbf{caminho de fallthrough}: código
tecnicamente morto quando toda saída da função termina com \texttt{RETURN},
mas necessário para que a função \texttt{main} (que não possui
\texttt{RETURN} explícito na IRT) retorne corretamente ao sistema operacional.

\section{Exemplo Completo: \texttt{factorial}}

\subsection{Fonte e IRT}

O programa TImp:

\begin{minted}{text}
function factorial(n : int) : int {
  if n == 0 then return 1
  else return n * factorial(n - 1)
}
function main() : void {
  print(factorial(read_int()))
}
\end{minted}

Gera (após canonicalização) a IRT mostrada no capítulo anterior.
Com a coloração \texttt{n}$\to$\texttt{\%rbx},
\texttt{t0..t2}$\to$\texttt{\%r12},
o compilador produz:

\subsection{Assembly Gerado (trecho de \texttt{factorial})}

\begin{minted}{asm}
.globl factorial
factorial:
    pushq  %rbp
    movq   %rsp, %rbp
    pushq  %rbx           # salva %rbx (callee-saved, usado por n)
    pushq  %r12           # salva %r12 (callee-saved, usado por t0/t1/t2)
    # loadParams: n (em %rdi) → %rbx
    movq   %rdi, %rbx
    # MOVE(t0, BINOP(EQ, n, 0)):  %rbx == 0?
    movq   %rbx, %rax
    pushq  %rax           # salva n
    movq   $0, %rax
    popq   %r11           # %r11 = n
    cmpq   %rax, %r11
    sete   %al
    movzbq %al, %rax
    movq   %rax, %r12     # t0 → %r12
    # CJUMP(t0, .true, .false)
    movq   %r12, %rax
    testq  %rax, %rax
    jne    .L_if_true_0
    jmp    .L_if_false_1
.L_if_true_0:
    movq   $1, %rax       # RETURN [1]
    popq   %r12
    popq   %rbx
    popq   %rbp
    retq
.L_if_false_1:
    # MOVE(t1, BINOP(SUB, n, 1))
    movq   %rbx, %rax
    pushq  %rax           # salva n
    movq   $1, %rax
    popq   %r11
    negq   %rax
    addq   %r11, %rax     # %rax = n - 1
    movq   %rax, %r12     # t1 → %r12
    # CALL(factorial, [t1]):  salva caller-saved antes
    movq   %r12, %rdi     # arg = t1
    xorq   %rax, %rax
    callq  factorial
    movq   %rax, %r12     # r → %r12
    # MOVE(t2, BINOP(MUL, n, r))
    movq   %rbx, %rax
    pushq  %rax           # salva n
    movq   %r12, %rax
    popq   %r11
    imulq  %r11, %rax     # %rax = n * r
    movq   %rax, %r12     # t2 → %r12
    # RETURN [t2]
    movq   %r12, %rax
    popq   %r12
    popq   %rbx
    popq   %rbp
    retq
.L_if_end_2:              # fallthrough (código morto)
    popq   %r12
    popq   %rbx
    popq   %rbp
    retq
\end{minted}

\begin{remark}
  Observe que os temporários \texttt{t0}, \texttt{t1} e \texttt{r} são
  todos mapeados a \texttt{\%r12}. Isso é correto porque os seus
  intervalos de vivacidade não se sobrepõem: \texttt{t0} morre antes de
  \texttt{t1} ser definido, e \texttt{t1} morre antes de \texttt{r} ser
  definido. A coloração por grafo captura exatamente essa compatibilidade.
\end{remark}

\section{O Pipeline Completo}

O back-end X86-64 é acionado pelo subcomando \texttt{--x86} do compilador
\texttt{timp}:

\begin{minted}{haskell}
run (Options CompileX86 fname) = do
  src <- readFile fname
  case timpParser src of
    Left err  -> putStrLn ("Parse error:\n" ++ err)
    Right ast ->
      case typeCheck ast of
        Left err -> putStrLn ("Type error: " ++ err)
        Right _  ->
          case compileTImp ast of
            Left err   -> putStrLn ("Code generation error: " ++ err)
            Right prog -> do
              let asmFile = fname -<.> "s"
              writeFile asmFile (pretty (irToX86 prog))
              putStrLn ("X86-64 assembly written to " ++ asmFile)
\end{minted}

O fluxo completo é:

\begin{center}
  \begin{tabular}{ll}
    \toprule
    Fase & Módulo principal \\
    \midrule
    Análise léxica        & \texttt{TImp.Frontend.Lexer.TImpLexer} \\
    Análise sintática     & \texttt{TImp.Frontend.Parser.TImpParser} \\
    Verificação de tipos  & \texttt{TImp.Frontend.TypeCheck.TImpTypeCheck} \\
    Geração de IRT        & \texttt{TImp.Backend.IR.TImpCodegen} \\
    Canonicalização       & \texttt{IR.Backend.X86.CFG}
    (\texttt{canonicalize}) \\
    Vivacidade            & \texttt{IR.Backend.X86.Liveness} \\
    Grafo de interferência & \texttt{IR.Backend.X86.IGraph} \\
    Intervalos            & \texttt{IR.Backend.X86.Intervals} \\
    Coloração             & \texttt{IR.Backend.X86.GraphColor} \\
    Seleção de instruções & \texttt{IR.Backend.X86.IRToX86} \\
    Impressão             & \texttt{IR.Backend.X86.X86Pretty} \\
    \bottomrule
  \end{tabular}
\end{center}

Para compilar, montar e executar:

\begin{verbatim}
cabal run timp -- --x86 --file examples/factorial.timp
gcc -o factorial examples/factorial.s -no-pie
echo 10 | ./factorial    # deve imprimir 3628800
\end{verbatim}

\section{Comparação dos Três Back-Ends}

O compilador TImp oferece três estratégias de geração de código, cada uma
com características distintas:

\begin{center}
  \begin{tabular}{llll}
    \toprule
    Aspecto & Back-end C & Back-end WAT & Back-end X86-64 \\
    \midrule
    Alocação de regs & compilador C & máquina de pilha & grafo de coloração \\
    ABI & do compilador C & do runtime WAT & System V AMD64 \\
    Portabilidade & qualquer plataforma & browser / WASI & Linux x86-64 \\
    Dependência externa & GCC / Clang & \texttt{wat2wasm} & apenas
    GCC (linker) \\
    Esforço de implem. & baixo & médio & alto \\
    Controle sobre o código & baixo & médio & total \\
    \bottomrule
  \end{tabular}
\end{center}

O back-end C delega o trabalho duro ao compilador C hospedeiro.
O back-end WAT tira proveito da máquina de pilha do WebAssembly para
simplificar a geração de código. O back-end X86-64 é o mais trabalhoso,
mas fecha o ciclo completo: nenhuma ferramenta de otimização externa é
necessária, e cada instrução gerada é resultado de uma decisão explícita
do compilador.

\section{Conclusão}

A geração de código X86-64 encerra o ciclo de compilação iniciado com a
análise léxica. Com a seleção de instruções concluída, o compilador TImp
é capaz de traduzir um programa de alto nível em instruções de máquina
reais, sem intermediários além do linker.

Os três capítulos deste módulo cobriram os desafios fundamentais de qualquer
back-end nativo:

\begin{itemize}
  \item \textbf{Canonicalização}: torna a IRT linear e adequada para análise
    de fluxo, eliminando efeitos colaterais escondidos em posições de expressão.
  \item \textbf{Vivacidade e interferência}: determinam quais temporários
    podem compartilhar registradores, formulando o problema de alocação como
    coloração de grafo.
  \item \textbf{Seleção de instruções}: traduz cada nó da IRT para instruções
    concretas, gerenciando o protocolo de chamada e o alinhamento de pilha.
\end{itemize}

As invariantes mantidas pelo gerador — resultado de expressão sempre em
\texttt{\%rax}, \texttt{\%rsp} múltiplo de 16 antes de todo \texttt{callq},
registradores callee-saved preservados em torno de cada retorno —
são as garantias que permitem ao código gerado interoperar corretamente
com a libc e com qualquer outro código compilado para o mesmo ABI.
